\scanpage{44}

\uline{Remarque} On voit facilement par des exemples
dans le Th et les deux corollaires, que
les conditions de platitude et de \uline{propreté}
sont essentielles. Cependant, il \deleted{semble} est
possible que l'hypothèse de quasi-projectivité de
$X$ \deleted{et celle de}\inserted{de $X'$}
puisse être omise, \deleted{et celle
de propreté \unsure{puisse} \unsure{omise} dans le Th, et les
\unsure{premiers cas}},
\unsure{sans doute non plus dans}
cor. 2 (ii) [Existence d'un schéma
algébrique \unsure{donnant} un schém. formel].

\margin{43}
